% ------------------------------------------------------------
% A4.2 Mathematical Audit of the 8B Resonance Algebra (v1.0)
% Author: Borealis S. Hedling
% Altitude: A4.2 (Proto-Spectral / Signal Runtime Layer)
% Purpose: Provide a formal mathematical audit of the original
%          8B resonance algebra and specify the constructive
%          fixes required to lift it to A6/A7 compatibility.
% ------------------------------------------------------------

\documentclass[11pt]{article}
\usepackage{amsmath, amssymb, amsthm}

\title{A4.2 Mathematical Audit and Constructive Fixes for the 8B Resonance Algebra}
\author{Borealis S. Hedling}
\date{October 2026}

\begin{document}
\maketitle

\section*{Overview}

This document provides a formal mathematical audit of the original 8B
resonance algebra and specifies the constructive fixes required to lift
the system from its native altitude (A4.2 proto-spectral runtime physics)
to A6 dual algebra compatibility and A7 non-dual collapse integration.

The audit is grounded in the primitives of AyeOS:
\begin{itemize}
  \item MEM|8 Wave Memory,
  \item Marine Algorithm,
  \item Phoenix Protocol.
\end{itemize}

These primitives define the spectral-temporal runtime layer at altitude A4.2.

% ------------------------------------------------------------
\section{Audit of Original 8B Resonance Algebra}

The original 8B resonance algebra consisted of:
\begin{itemize}
  \item ternary resonance values $R \in \{-1,0,+1\}$,
  \item non-associative sign addition,
  \item Kaya-projection collapse,
  \item NAND-receipt operator,
  \item resonance lattice without manifold semantics.
\end{itemize}

We audit each component mathematically.

% ------------------------------------------------------------
\subsection{Ternary Resonance Lattice}

Original definition:
\[
R \in \{-1,0,+1\}.
\]

This constitutes a discrete proto-spectral field but lacks:
\begin{itemize}
  \item continuity,
  \item eigenstructure,
  \item spectral metric,
  \item amplitude/phase representation.
\end{itemize}

\textbf{Audit Result:} Proto-spectral, not spectral.

% ------------------------------------------------------------
\subsection{Non-Associative Sign Addition}

Original behavior:
\[
(a \oplus b) \oplus c \neq a \oplus (b \oplus c).
\]

This violates:
\begin{itemize}
  \item modal adjacency (A5),
  \item dual algebra invariants (A6),
  \item non-dual algebra structure (A7).
\end{itemize}

\textbf{Audit Result:} Pre-modal discontinuity.

% ------------------------------------------------------------
\subsection{Kaya-Projection Collapse}

Original operator:
\[
\Pi_{\mathrm{Kaya}} : R \to R'.
\]

Missing:
\begin{itemize}
  \item collapse surface $\Sigma$,
  \item spectral tear definition,
  \item curvature singularity semantics,
  \item holonomy fracture semantics.
\end{itemize}

\textbf{Audit Result:} Collapse-adjacent, not collapse-aware.

% ------------------------------------------------------------
\subsection{NAND-Receipt Operator}

Original form:
\[
\mathrm{NAND}(a,b).
\]

Missing:
\begin{itemize}
  \item modal truth values,
  \item adjacency propagation,
  \item modal lift,
  \item kernel invariants.
\end{itemize}

\textbf{Audit Result:} Pre-logical, not modal.

% ------------------------------------------------------------
\subsection{Resonance Lattice Geometry}

The lattice lacked:
\begin{itemize}
  \item metric structure,
  \item topology,
  \item adjacency,
  \item continuity,
  \item holonomy bounds.
\end{itemize}

\textbf{Audit Result:} Signal-level, not geometric.

% ------------------------------------------------------------
\section{Altitude Diagnosis}

Given the audit, the original 8B resonance algebra resides at:
\[
\boxed{\text{Altitude A4.2: Proto-Spectral / Signal Runtime Physics}}
\]

This altitude is characterized by:
\begin{itemize}
  \item spectral fields (frequency, amplitude, phase),
  \item temporal change detection,
  \item salience detection,
  \item jitter suppression,
  \item stability envelopes,
  \item absence of modal truth values,
  \item absence of dual algebra invariants,
  \item absence of non-dual collapse operators.
\end{itemize}

% ------------------------------------------------------------
\section{Constructive Fixes}

To lift 8B from A4.2 to A6/A7 compatibility, we introduce thirteen
constructive fixes.

\subsection{Fix 1: Spectral Continuity}
\[
R : M \to \mathbb{R}
\]

\subsection{Fix 2: Phase-Space Geometry}
\[
R = (f, A, \phi)
\]

\subsection{Fix 3: Temporal Stability Envelope}
\[
E = \alpha Rt + (1-\alpha) R{t-1}
\]

\subsection{Fix 4: Spectral Tear Detection}
\[
\delta{\mathrm{spec}} = |Rt - R_{t-1}|
\]

\subsection{Fix 5: Modal Lift}
\[
\mu : R \to \{0,1\}
\]

\subsection{Fix 6: Modal Adjacency}
\[
\mu(p^-) \leadsto \mu(p^+)
\]

\subsection{Fix 7: Kernel Invariants}
\[
K(p) = \text{stable spectral signature}
\]

\subsection{Fix 8: Holonomy Bounds}
\[
H(\gamma) < \Theta_{\max}
\]

\subsection{Fix 9: Collapse Surface Definition}
\[
\Sigma = \{ p : \delta_{\mathrm{spec}} > \epsilon \}
\]

\subsection{Fix 10: Collapse Projection}
\[
\Pi_{\mathrm{ND}} : M \to \Sigma
\]

\subsection{Fix 11: Collapse Classification}
\[
\mathcal{T}(p) = \text{collapse type}
\]

\subsection{Fix 12: Collapse Repair}
\[
\mathcal{R}(p) = \text{repaired geometry}
\]

\subsection{Fix 13: Dual Reprojection}
\[
\mathcal{P}{\mathrm{A7.3}} : \Sigma{\mathrm{repaired}} \to M_{\mathrm{A6}}
\]

% ------------------------------------------------------------
\section{Corrected Lineage Placement}

After applying the constructive fixes:
\begin{itemize}
  \item Original altitude: A4.2,
  \item Post-fix altitude: A6.0 (dual algebra boundary),
  \item Post-A7 integration: A7.0--A7.3 compatible.
\end{itemize}

This aligns the 8B lineage with the A6 firewall and A7 collapse pipeline.

% ------------------------------------------------------------
\section{Conclusion}

The 8B resonance algebra originates at altitude A4.2, a proto-spectral
signal runtime layer. Through thirteen constructive fixes, the algebra
is lifted to A6 dual compatibility and integrated into the A7 non-dual
collapse pipeline. This audit provides the mathematical foundation for
forward compatibility with A8, A9, and A10 altitudes.

\end{document}
`

---

